\subsection*{The Missing Identity}

\begin{remark*}[Exposition]
In the one-based system, every element of $\mathbb{N}$ is positive. The
operation $+$ is defined and behaves associatively and commutatively, but
there is no element $e\in\mathbb{N}$ satisfying $e+n=n$ for every
$n\in\mathbb{N}$. The element needed for that identity is not another
successor-stage element; it is a new element placed before $1$.
\end{remark*}

\begin{remark*}[Interpretation]
The passage from $\mathbb{N}$ to $\mathbb{W}$ is therefore not a change in the
positive arithmetic already constructed. It is an extension that adds the
missing additive identity and leaves the positive part intact.
\end{remark*}

\begin{dependencies}
\begin{itemize}
  \item \hyperref[def:addition-on-peano-system]{Addition on a Peano System}
  \item \hyperref[thm:addition-associative]{Addition Is Associative}
  \item \hyperref[thm:addition-commutative]{Addition Is Commutative}
\end{itemize}
\end{dependencies}

\begin{remark*}[Exposition]
Once $0$ is available, addition can be packaged as a genuine commutative
monoid. That structure is the natural input for the integer construction,
where ordered pairs are interpreted as formal differences and subtraction is
made available by a quotient construction.
\end{remark*}
